\documentclass[conference]{IEEEtran}
\IEEEoverridecommandlockouts
\usepackage{cite}
\usepackage{amsmath,amssymb,amsfonts}
\usepackage{algorithmic}
\usepackage{graphicx}
\usepackage{textcomp}
\usepackage{xcolor}
\usepackage{hyperref}
\usepackage{listings}
\usepackage{tabularx}
\usepackage{balance}

\newcolumntype{C}{>{\centering\arraybackslash}X}

\usepackage{tikz}
\usetikzlibrary{shapes,arrows,positioning,fit,backgrounds}

\def\BibTeX{{\rm B\kern-.05em{\sc i\kern-.025em b}\kern-.08em
    T\kern-.1667em\lower.7ex\hbox{E}\kern-.125emX}}

\lstset{
  basicstyle=\ttfamily\scriptsize,
  columns=fullflexible,
  frame=single,
  breaklines=true,
  captionpos=b,
}

\begin{document}

\title{Post-Quantum Workload Identity and Long-Term Audit Integrity for MCP-Based AI Agents}

\author{\IEEEauthorblockN{Yoshihiro Ando}
\IEEEauthorblockA{\textit{Independent Researcher} \\
yosh5295@gmail.com \\
Tokyo, Japan}
}

\maketitle

\begin{abstract}
As autonomous AI agents transition into a globally interconnected ecosystem via the Model Context Protocol (MCP), securing their identity and the long-term confidentiality of their reasoning history against quantum-era threats is a prerequisite for sovereign autonomy. This paper, representing \textbf{Phase 3: Post-Quantum (Temporal Security)} of a comprehensive security trilogy, proposes a Post-Quantum Cryptographic (PQC) framework. We introduce the \textit{PQC-Audit-Chain}, which ensures non-repudiable integrity through NIST-standardized ML-DSA signatures and long-term confidentiality via ML-KEM hybrid encryption. We present \textbf{PQC Remote Attestation} for client-side integrity proof and \textbf{Context-Adaptive Security Scaling (CASS)} for risk-aware agility. By establishing \textit{Temporal Integrity Proofs}, we ensure that an agent's history remains immune to retroactive tampering and decryption by future quantum adversaries.
\end{abstract}

\begin{IEEEkeywords}
Post-Quantum Cryptography, AI Agents, Model Context Protocol, ML-DSA, ML-KEM, Remote Attestation, Hybrid Encryption, Temporal Integrity.
\end{IEEEkeywords}

\section{Introduction}
The rapid advancement of LLM agents has created a dependency on long-lived audit logs and identity tokens. Phases 1 and 2 \cite{ando2026guardrails, ando2026zerotrust} established structural and behavioral security. However, these systems rely on classical cryptography (RSA/ECDSA), which is vulnerable to Shor's algorithm on future quantum computers.

A primary threat in the agentic era is **"Store Now, Decrypt Later" (SNDL)**. Adversaries can capture current encrypted agent traffic (tool calls, reasoning traces, proprietary data) and store it until a Cryptographically Relevant Quantum Computer (CRQC) is available to decrypt it.

This paper concludes our technical trilogy:
\begin{itemize}
    \item \textbf{Phase 1: Guardrails (Pre-execution Security) \cite{ando2026guardrails}.} Structural containment.
    \item \textbf{Phase 2: Zero Trust (Execution Security) \cite{ando2026zerotrust}.} Identity and behavioral IDS (Mamba).
    \item \textbf{Phase 3: Post-Quantum (Temporal Security).} The current work, focusing on \textbf{Temporal Integrity} and future-proof identities.
\end{itemize}

\section{Post-Quantum Hybrid Resilience}
To ensure a smooth transition and maintain compatibility, \texttt{llm-cli} implements a hybrid approach combining classical and post-quantum primitives.

\subsection{Quantum-Safe Identity (ML-DSA)}
We upgrade the workload identity using NIST FIPS 204 (ML-DSA) \cite{nistfips204}. Every tool execution request is signed with a hybrid token:
\begin{equation}
\sigma_{hybrid} = (\sigma_{RSA2048}, \sigma_{ML-DSA-65})
\end{equation}
ML-DSA-65 (Module-Lattice-based Digital Signature Algorithm) provides security equivalent to AES-192. This dual-signature approach ensures that even if one primitive is compromised (by quantum or classical means), the non-repudiability of the agent's action remains intact.

\subsection{Confidential Audit Logging (ML-KEM)}
Audit logs contain sensitive reasoning traces, tool arguments (including potential PII), and system responses. We implement \textbf{Hybrid-PQC Encryption} using NIST FIPS 203 (ML-KEM) \cite{nistfips203} and AES-256-GCM. 
The system encapsulates an AES-256 key using ML-KEM-768. This key is then used to encrypt the audit log payload. By protecting the "Temporal Dimension" of agent data, we ensure that an agent's history cannot be retroactively parsed by future quantum adversaries to extract trade secrets or intellectual property.

\section{PQC Remote Attestation}
In a distributed environment, a server must trust that the client code has not been tampered with to bypass security guardrails. We introduce \textbf{PQC Remote Attestation}:
\begin{enumerate}
    \item \textbf{Manifest Generation:} The client generates a SHA-256 hash manifest of its core source code and configuration.
    \item \textbf{Signing:} This manifest is signed using the client's ML-DSA private key.
    \item \textbf{Attestation Token:} The resulting token is embedded in the Hybrid Identity JWT.
\end{enumerate}
This allows remote MCP servers to verify that the agent is running on an "Authentic Stack" where Phase 1 guardrails are guaranteed to be active.

\section{Context-Adaptive Security Scaling (CASS)}
The overhead of PQC (larger signatures and keys) can impact performance. We propose \textbf{CASS}, a mechanism that dynamically scales security levels based on the tool's risk profile:
\begin{itemize}
    \item \textbf{Low Risk} (e.g., \texttt{search\_web}): Uses classical RSA and standard monitoring to minimize latency.
    \item \textbf{High Risk} (e.g., \texttt{execute\_python}, \texttt{edit\_file}): Escalates to full Hybrid-PQC signatures, KEM-encrypted logging, and maximum Mamba Sentinel sensitivity.
\end{itemize}
CASS ensures "Security with Agility," applying expensive quantum-resistant measures only when the potential impact of a compromise justifies the cost.

\begin{figure*}[htbp]
\centering
\resizebox{0.85\textwidth}{!}{
    \input{trilogy\_phase3_architecture.tex}
}
\caption{Comprehensive Security Trilogy Architecture: Integrating Phase 2 (Mamba Sentinel) and Phase 3 (PQC Hybrid Identity) for sovereign AI governance.}
\label{fig:pqc_arch}
\end{figure*}

\section{Evaluation}
Benchmarks were performed on an \textbf{AMD Ryzen 5 8540U} system.

\subsection{Cryptographic Performance (Empirical)}
Table \ref{tab:pqc_bench} shows the latencies of NIST PQC primitives.

\begin{table}[htbp]
\caption{ML-DSA Performance (Latency in ms)}
\begin{center}
\begin{tabularx}{\linewidth}{|l|C|C|}
\hline
\textbf{Algorithm} & \textbf{Sign (ms)} & \textbf{Verify (ms)} \\
\hline
ML-DSA-65 & 23.30 & 6.58 \\
\hline
\end{tabularx}
\label{tab:pqc_bench}
\end{center}
\end{table}

\subsection{The Cost of Future-Proofing}
While ML-DSA-65 introduces a ~23ms signing latency, this overhead is less than 2\% of a typical LLM reasoning cycle (usually >1,000ms). The trade-off—gaining immunity against future quantum decryption and current spoofing—is highly favorable for high-stakes autonomous systems.

\section{Conclusion}
This completion of our technical trilogy establishes a foundational roadmap for sovereign AI agents. By combining structural guardrails (Phase 1), behavioral IDS (Phase 2), and post-quantum integrity (Phase 3), we ensure that autonomous systems can operate safely across spatial, behavioral, and temporal dimensions. As agents become the primary interface for system interaction, these "Triple-Lock" security standards will be essential for maintaining human control and data sovereignty in the AI era.

\begin{thebibliography}{00}
\bibitem{mcp2024} Anthropic, ``Model Context Protocol (MCP) Specification,'' 2024.
\bibitem{nistfips203} NIST, ``FIPS 203: ML-KEM Standard,'' 2024.
\bibitem{nistfips204} NIST, ``FIPS 204: ML-DSA Standard,'' 2024.
\bibitem{ando2026guardrails} Y. Ando, ``Secure-by-Default Guardrails for MCP-Based Tool Use in Multi-Modal LLM Agents,'' \textit{TechRxiv}, 2026.
\bibitem{ando2026zerotrust} Y. Ando, ``Zero-Trust Architecture for MCP-Based AI Agents,'' \textit{TechRxiv}, 2026.
\end{thebibliography}

\end{document}
